\documentclass{situationist-preprint}

\title{Dirty Hands and Appropriate Remorse}
\msid{ETH-27-a}
\msfamily{ETH-27}
\msdiscipline{Normative and Applied Ethics}
\mstype{Principal manuscript}
\mskeywords{dirty hands; remorse; moral remainder; agent-regret; infringement and violation; fitting attitudes; political ethics}

\begin{document}
\maketitle

\begin{abstract}
A politician who can avert a grave evil only by deceiving, bargaining with the corrupt, or breaking a solemn undertaking is said to dirty her hands. Walzer argued that she does wrong in doing right: her act is justified, she is guilty, and we should want her to feel it. The view is unstable, since it requires a single act to be both what she ought to do and something she ought not to do. We argue that the agent who dirties her hands for a greater good acts rightly, and that remorse remains appropriate. The key move is to separate two evaluations that Walzer runs together. Whether the act is right is settled by the balance of reasons; whether remorse is fitting is settled by what remorse represents, namely that one has, by one's own agency, infringed a serious claim of another. Dirty-hands acts are right and infringe such claims, so remorse fits them, while the judgement that one acted wrongly does not. The account explains the duties of acknowledgment and amends that survive the act, and it distinguishes remorse both from guilt and from Williams's agent-regret.
\end{abstract}

\section{The problem}

Machiavelli advised the prince to learn how not to be good, and to use that knowledge according to necessity \citep[ch.~15]{Machiavelli1532}. \citet{Weber1919} put the same thought in the vocabulary of the ethics of responsibility: whoever goes into politics makes a pact with diabolical powers, and cannot count on good following only from good. The modern problem of dirty hands was formulated by \citet{Walzer1973}, who took his title from Sartre's play \citep{Sartre1948}. Walzer's cases are two. A candidate who wants to win an election in order to do good must make a corrupt deal with a ward boss. A leader who has long opposed torture, and who has won office pledged to end a colonial war, must decide whether to order the torture of a captured rebel who knows where bombs have been hidden that will soon explode in apartment buildings. Walzer's claim is that in such cases the politician does the right thing, that in doing it she does something wrong, and that she is guilty and ought to feel guilty. A politician who felt nothing, Walzer thought, would show that she was not a good person, however successful she might be.

Walzer's view has an obvious difficulty. If the act is right, it is what the agent ought to do all things considered; if it is wrong, it is something she ought not to do all things considered. A single act cannot be both, unless there are genuine moral dilemmas in which every option is wrong, and Walzer's case would be one, since refraining from the act is also wrong. ETH-15 argues that there are no genuine dilemmas: every apparent dilemma has a uniquely permissible option. This leaves two other ways of understanding dirty hands. On the first, which we may call the \emph{clean-hands} view, the right act is simply right, and the residue of feeling is a psychological fact with no normative standing, useful perhaps as a sign of a good disposition \citep{Hare1981}. On the second, the agent's remorse is appropriate to something that her act really involved, though not to its being wrong. We defend the second. The agent who dirties her hands for a greater good acts rightly, and remorse remains appropriate.

The difficulty for the second view is to say what remorse is appropriate to, if not to wrongdoing. \citet{Williams1976} introduced \emph{agent-regret}, the regret an agent feels about her own past actions, including their unintended and involuntary aspects, as distinct from the regret a spectator may feel about a state of affairs; his lorry driver who, through no fault of his own, runs over a child, feels something that no spectator can feel. Agent-regret is appropriate without wrongdoing, but it is too weak for our case. The lorry driver did nothing he chose to do to the child, whereas the politician knowingly deceived, bargained or broke her word. What is needed is an object for remorse that lies between wrongdoing, which would make the case a dilemma, and mere involvement, which would assimilate it to the driver's.

\section{Framework}

\begin{definition}[Infringement and violation]\label{d:infringe}
A person has a \emph{claim} against an agent when the agent owes it to that person to act or refrain in some way. The agent \emph{infringes} the claim when she fails so to act; she \emph{violates} it when she infringes it and her infringing it is impermissible all things considered \citep[ch.~3]{Thomson1990}.
\end{definition}

\begin{definition}[Dirty-hands case]\label{d:dirty}
A \emph{dirty-hands case} is a situation in which an agent can secure a great good, or avert a great evil, only by an act that infringes a serious claim of some person, and in which the good secured or evil averted outweighs, on the balance of reasons, the claim infringed.
\end{definition}

Definition~\ref{d:dirty} builds the outweighing into the case. Whether Walzer's torture case, or any other actual case, satisfies it is a first-order question that we do not settle; the thesis concerns the cases that satisfy it, and Walzer's thesis concerns the same cases, since he holds that the act is justified.

\begin{definition}[Remorse]\label{d:remorse}
\emph{Remorse} is a painful, backward-looking emotion directed at one's own past act, which presents that act as having infringed, through one's own agency, a serious claim of another, and which disposes the agent to acknowledge the infringement to the person whose claim it was, to make amends, and to be reluctant to act in the same way again.
\end{definition}

\begin{principle}[Fittingness]\label{p:fitting}
An emotion is fitting if and only if what it presents its object as being is the case.
\end{principle}

\begin{assumption}[No dilemmas]\label{a:etb15}
In every situation in which each option infringes some claim, exactly one option is permissible all things considered, and it is the option best supported by the balance of reasons.
\end{assumption}

Principle~\ref{p:fitting} is the fitting-attitude conception of the appropriateness of emotions. It separates the question whether an emotion is fitting from the question whether it is good, prudent or morally admirable to feel it, a separation whose neglect \citet{DArmsJacobson2000} call the moralistic fallacy. It presupposes that emotions present their objects as having evaluative properties, which MND-17 defends: fear presents danger, and on Definition~\ref{d:remorse} remorse presents infringement. Assumption~\ref{a:etb15} is the conclusion of ETH-15, together with the balance principle on which it rests.

Definition~\ref{d:remorse} is the substantive premise. It distinguishes remorse from two neighbors. \emph{Guilt}, in the sense Walzer intends, involves the judgement that one acted wrongly, and so presents one's act as a violation. \emph{Agent-regret}, in Williams's sense, presents one's act only as one's own and as having had a bad upshot. Remorse is between them: it presents the act as an infringement of a claim, which entails that someone was owed something and did not get it, but not that the agent should have done otherwise. The definition is supported by the dispositions that characterize remorse and that distinguish it from agent-regret. Remorse is directed at a person, the one whose claim was infringed, and it seeks acknowledgment and repair from the agent to that person; agent-regret may seek repair, but it is not addressed to anyone as the bearer of a claim. And remorse is extinguished by learning that the person had no claim, as when the agent discovers that the person she deceived had released her in advance from any duty of candor, whereas it is not extinguished by learning that the act was justified.

\section{The argument}

\begin{theorem}\label{t:right}
In a dirty-hands case, the agent who performs the dirty-hands act acts rightly.
\end{theorem}

\begin{proof}
In a dirty-hands case, refraining from the act fails to secure the good or avert the evil, and in the cases that interest us this too infringes claims, those of the people the evil will befall. By Definition~\ref{d:dirty}, the dirty-hands act is the option best supported by the balance of reasons. By Assumption~\ref{a:etb15}, it is the uniquely permissible option. To perform the uniquely permissible option is to act rightly.
\end{proof}

\begin{theorem}\label{t:remorse}
In a dirty-hands case, remorse for the dirty-hands act is fitting.
\end{theorem}

\begin{proof}
By Definition~\ref{d:remorse}, remorse presents the act as having infringed, through the agent's own agency, a serious claim of another. By Definition~\ref{d:dirty}, the act infringes a serious claim of some person, and the agent performs it. What remorse presents is therefore the case, and by Principle~\ref{p:fitting} remorse is fitting.
\end{proof}

\begin{corollary}\label{c:guilt}
In a dirty-hands case, guilt, understood as presenting the act as wrong, is not fitting.
\end{corollary}

\begin{proof}
By Theorem~\ref{t:right}, the act is permissible, so the claim it infringes is not violated (Definition~\ref{d:infringe}). Guilt in the relevant sense presents the act as a violation, and by Principle~\ref{p:fitting} it is not fitting.
\end{proof}

Theorems~\ref{t:right} and~\ref{t:remorse} together are the thesis. They are consistent because they answer different questions with different standards: the first asks whether the act was the one to perform, and answers by the balance of reasons; the second asks whether a certain evaluative presentation of the act is correct, and answers by whether the act had the property presented. The infringed claim figures in both. In the first it is outweighed, and outweighing settles what to do. In the second it is present, and presence is all that the content of remorse requires. That an outweighed claim persists, and grounds further duties, is the point of Corollary~3.5 of ETH-15-a, which shows that in every apparent dilemma the permissible option infringes a requirement in force.

\begin{proposition}\label{p:amends}
In a dirty-hands case, the agent who acts rightly owes the person whose claim she infringed acknowledgment and, where possible, amends.
\end{proposition}

\begin{proof}
A claim is something owed to a person. When it is permissibly infringed, the infringement is justified, but the person has still not received what she was owed. Feinberg's example makes the point: the backpacker who breaks into a cabin to survive a blizzard acts permissibly, and owes the owner compensation for what he damages and uses \citep{Feinberg1978}; Thomson generalizes it into the thesis that permissible infringements of claims leave residual duties \citep[ch.~3]{Thomson1990}. The debt is not discharged by the justification of the infringement, since the justification concerns the agent's choice, while the debt concerns the person's loss. The dispositions of remorse in Definition~\ref{d:remorse} are dispositions to pay this debt, and their fittingness is of a piece with the debt's existence.
\end{proof}

Proposition~\ref{p:amends} gives Walzer's best observation its proper form. Walzer suggested that the politician who dirties her hands should be willing to bear some public cost, and drew on Camus's just assassins, who accept their own deaths, as a model. On our account the cost is not punishment, which presupposes wrongdoing, but acknowledgment and amends to those whose claims were infringed. The same structure appears where an agent breaks a just law to prevent grave harm: ETH-17 argues that this is permissible only when the breach is public and its legal consequences are accepted, and the acceptance plays the role that amends play here.

\section{Cases}

\paragraph{The broken undertaking.} A minister has promised a community that a factory will not close. Circumstances change, and keeping the factory open would require diverting funds from a program that will prevent many deaths. She diverts the funds. Her act is right, and the community's claim, grounded in the promise, has been infringed. She owes them an explanation in person, not by press release, and such help in the transition as is in her power. A minister who sends the press release and feels nothing has misperceived her situation, even though she has chosen correctly.

\paragraph{The trolley bystander.} A bystander diverts a trolley from five people onto a track where it kills one. ETH-03 argues that this is permissible, because the one has no claim that an existing threat not be redirected towards him, provided it is redirected so as to save a greater number. On the present account, then, the bystander's act infringes no claim, and remorse is not fitting. Agent-regret is: the bystander's act caused a death, and she can rightly wish that she had not been the one through whom it came. The contrast confirms Definition~\ref{d:remorse}. Remorse tracks the infringement of claims, and where the best moral theory of the case finds none, remorse is out of place even though a death was caused.

\paragraph{The lorry driver.} Williams's driver who, without fault, runs over a child, infringes no claim either, on the plausible view that the child's claim against drivers is a claim not to be harmed through their fault. Agent-regret is fitting, remorse is not, and the driver who nonetheless feels something like remorse feels it for the reason Williams gives: an agent who could simply shrug off his causal involvement in a death would be lacking in something. ETH-13 discusses the related case of the careless driver in connection with moral luck.

\section{Objections}

\begin{objection}[The name]
If the agent acts rightly and violates no claim, her hands are clean. The problem of dirty hands has been dissolved, not solved.
\end{objection}

\begin{reply}
The dirt is the infringement. A claim is something owed to a person, and the agent has deliberately failed to give a person what was owed to him. That is a moral fact about her act, and it remains true whatever the act's justification. What has been dissolved is only the paradox that the act is both what she ought to do and what she ought not to do. Walzer needed the paradox to explain why she should feel guilty; we explain why she should feel remorse without it.
\end{reply}

\begin{objection}[Irrationality]
An emotion that is painful, that presents its object as regrettable, and that disposes one not to do the same again, is irrational when directed at an act one would, on reflection, perform again. The clean-hands view is right that such remorse is at best a useful disposition \citep{Hare1981}.
\end{objection}

\begin{reply}
The objection assumes that remorse presents the act as one that should not have been done. On Definition~\ref{d:remorse} it does not. An agent who would make the same choice again and feels remorse is not divided against herself: she judges that she should have done it, and she perceives that someone was wronged in the sense of having a claim infringed. The disposition to reluctance is not a disposition to refuse; it is a disposition to require that the case be clear before acting, which is the disposition Walzer rightly wanted politicians to have. The clean-hands view must say that the minister who sends the press release and feels nothing has made no error. We think she has made an error of perception, and Theorem~\ref{t:remorse} identifies it.
\end{reply}

\begin{objection}[Inescapable wrongdoing]
\citet{Gowans1994} argues that in conflicts of this kind the agent does wrong to persons even if she does not do the wrong thing, and that only the recognition of inescapable wrongdoing makes sense of the moral emotions involved. Remorse responds to wronging, not to infringing.
\end{objection}

\begin{reply}
We accept the distinction between wronging a person and doing the wrong thing, and our account draws it in Thomson's terms: to infringe a claim is to fail to give a person what is owed to him, which is one sense of wronging him; to violate it is to do so impermissibly. If Gowans's \emph{wronging} is infringement, our view and his coincide on the emotions. If it is violation, then the agent who wrongs someone in his sense acts impermissibly, and whenever all her options would wrong someone, she faces a genuine dilemma, which ETH-15 excludes. Gowans's phenomenology is therefore best preserved by our reading.
\end{reply}

\begin{objection}[The admirable agent]
\citet{Greenspan1995} argues that the guilt feelings that persist after moral conflict are not responses to the agent's having done wrong, but are warranted as expressions of the agent's commitment to the norms she has been forced to override. If so, the appropriate emotion is grounded in the agent's character, not in a feature of the act, and Definition~\ref{d:remorse} misplaces it.
\end{objection}

\begin{reply}
The two accounts are complementary. That a good agent will feel remorse, and that her feeling it expresses her commitment to the norms she overrode, is a claim about the value of the emotion and about what its presence reveals. Principle~\ref{p:fitting} concerns its fittingness, which depends on its object. An emotion can be both fitting and an expression of good character, and remorse in dirty-hands cases is both. What Greenspan's account cannot do alone is explain why the remorse should be directed at the particular person whose claim was overridden, and why it should seek repair from that person, which is what Definition~\ref{d:remorse} explains.
\end{reply}

\section{What the result does not establish}

The thesis concerns the cases that satisfy Definition~\ref{d:dirty}, in which the good outweighs the infringed claim. It does not say that any given act of political deception, bargaining or coercion satisfies it, and it does not say that torture ever does: whether some claims are so stringent that no good can outweigh them is a first-order question that the definition leaves open, and on which the thesis is neutral. Nor does the thesis say that remorse is required of the agent in the sense that its absence is itself a further wrong; it says that remorse is fitting, and that the duties of acknowledgment and amends of Proposition~\ref{p:amends} are owed. The companion manuscript ETH-27-b develops the account of remorse without wrongdoing outside political contexts.

\section{Conclusion}

The agent who dirties her hands for a greater good acts rightly, since her act is the uniquely permissible option in a situation in which every option infringes some claim. Remorse remains appropriate, because remorse presents an act as the infringement of a serious claim of another, and her act was such an infringement. Guilt, as the judgement that she acted wrongly, is not fitting, and agent-regret is too weak to capture what she owes. The residue of the dirty-hands case is a debt to the person whose claim was overridden, and remorse is the emotion that registers it.

\begin{thebibliography}{99}\small
\bibitem[D'Arms and Jacobson(2000)]{DArmsJacobson2000} D'Arms, J. and Jacobson, D. (2000). The moralistic fallacy: On the `appropriateness' of emotions. \emph{Philosophy and Phenomenological Research}, 61(1), 65--90.
\bibitem[Feinberg(1978)]{Feinberg1978} Feinberg, J. (1978). Voluntary euthanasia and the inalienable right to life. \emph{Philosophy \& Public Affairs}, 7(2), 93--123.
\bibitem[Gowans(1994)]{Gowans1994} Gowans, C.~W. (1994). \emph{Innocence Lost: An Examination of Inescapable Moral Wrongdoing}. Oxford University Press.
\bibitem[Greenspan(1995)]{Greenspan1995} Greenspan, P.~S. (1995). \emph{Practical Guilt: Moral Dilemmas, Emotions, and Social Norms}. Oxford University Press.
\bibitem[Hare(1981)]{Hare1981} Hare, R.~M. (1981). \emph{Moral Thinking: Its Levels, Method, and Point}. Oxford University Press.
\bibitem[Machiavelli(1532)]{Machiavelli1532} Machiavelli, N. (1532). \emph{Il Principe}. Rome: Antonio Blado. Cited by chapter.
\bibitem[Sartre(1948)]{Sartre1948} Sartre, J.-P. (1948). \emph{Les Mains sales}. Paris: Gallimard.
\bibitem[Thomson(1990)]{Thomson1990} Thomson, J.~J. (1990). \emph{The Realm of Rights}. Harvard University Press.
\bibitem[Walzer(1973)]{Walzer1973} Walzer, M. (1973). Political action: The problem of dirty hands. \emph{Philosophy \& Public Affairs}, 2(2), 160--180.
\bibitem[Weber(1919)]{Weber1919} Weber, M. (1919). \emph{Politik als Beruf}. Munich: Duncker \& Humblot.
\bibitem[Williams(1976)]{Williams1976} Williams, B. (1976). Moral luck. \emph{Proceedings of the Aristotelian Society, Supplementary Volumes}, 50, 115--135.
\end{thebibliography}

\end{document}
